% !TeX program = latexmk
% Full report: cover, abstract, contents, symbols and abbreviations, body,
% references, appendix.
% Build: make (or latexmk -pdf main.tex). Style: tuereport.sty.
% References: references.bib. Writing rules: WRITING.md in the LaTeX templates.
\documentclass[a4paper, 11pt]{article}
\usepackage{tuereport}
\addbibresource{references.bib}

\reporttitle{Title of the report}
\reportsubtitle{COURSE CODE -- Course name\\Quartile, academic year}
\reportversion{v0.1}                      % v0.x drafts, v1.0 first release
\reportgroup{Group XX}                    % leave empty ({}) to hide
\reportauthors{%
  Firstname Lastname & 1234567\\
  Firstname Lastname & 1234567\\
}
\reportsupervisor{Tutor: Name Surname}
\reportkeywords{}                          % PDF keywords, comma-separated
% \reportcoverimage{figures/cover.png}    % optional picture on the cover
% \reportdate{25 November 2026}            % default: \today
% \accentcolor{HTML}{003399}              % secondary color for headings etc.; default Scarlet

\begin{document}

\makecover

\frontmatter
% Abstract: the micro version of the report. Trends leading to the problem, the
% problem, what is new (approach, solution), and the main results. No citations.
\section*{Abstract}
\addcontentsline{toc}{section}{Abstract}
State the problem in one or two sentences.
Say what was done to solve it and what is new about the approach.
Give the main result with numbers, for example a launch distance of \SI{9.3}{\metre}.

\tableofcontents
\clearpage

\section*{List of symbols}
\addcontentsline{toc}{section}{List of symbols}
\label{sec:symbols}
\begin{table}[H]
  \rowcolors{2}{gray!10}{white}
  \centering
  \begin{tabular}{m{0.14\textwidth} m{0.58\textwidth} m{0.15\textwidth}}
    \toprule
    \textbf{Symbol} & \textbf{Meaning} & \textbf{Unit} \\
    \midrule
    $F$ & Force exerted by the spring & \si{\newton} \\
    $k$ & Spring constant & \si{\newton\per\metre} \\
    $e_\txt{total}$ & Total efficiency & -- \\
    \bottomrule
  \end{tabular}
\end{table}

\section*{List of abbreviations}
\addcontentsline{toc}{section}{List of abbreviations}
% Only abbreviations used with \ac{...} are printed (with the page of first
% use).
\begin{acronym}[SOC]
  \acro{SOC}{system on chip}
  \acro{ADC}{analog-to-digital converter}
\end{acronym}

% Optional for long reports:
% \listoffigures
% \listoftables

\mainmatter
% ------------------------------------------------------------------------------
\section{Introduction}
\label{sec:introduction}
% The mini version of the report: trends leading to the problem, the problem
% statement, what is new, the approach and solution, and how the report is
% organized.
Start with the context that leads to the problem.
State the problem precisely and the requirements a solution must meet.
Summarize the approach and the main claims, which the rest of the report proves.

The report is organized as follows.
\Cref{sec:related} compares this work with existing solutions.
\Cref{sec:approach} describes the approach, \cref{sec:setup} the experimental setup, and \cref{sec:results} the results.
\Cref{sec:conclusions} concludes.

% ------------------------------------------------------------------------------
\section{Related work}
\label{sec:related}
% How this work differs from, improves on, or extends other work. Not just a
% list.
Earlier designs used a counterweight~\cite{knuth1986} or a spring~\cite{bbc_springs}.
Unlike both, this work \todo{state the difference}.

% ------------------------------------------------------------------------------
\section{Approach}
\label{sec:approach}
% Short introduction linking to the previous section, subsections, short
% conclusion linking to the next section.
This section models the mechanism and derives the design parameters.

\subsection{Model}
\label{sec:model}
The spring force follows Hooke's law, given in \cref{eq:hooke}.
\begin{equation}
  F = k \, \delta
  \label{eq:hooke}
\end{equation}
With $k = \SI{250}{\newton\per\metre}$ and a compression of \SI{40}{\milli\metre}, the force is \SI{10}{\newton}.
\begin{block}{Claim}
  The design reaches \SI{10}{\metre} with \SI{20}{\gram}.
\end{block}
The \ac{ADC} on the \ac{SOC} samples the force sensor; a second use of \ac{ADC} is not expanded again.

\subsection{Design}
\label{sec:design}
\Cref{fig:overview} shows the parts of the design, and \cref{fig:steps} the three stages of a launch.
Explain what the reader should learn from each figure, not only what it shows.

\begin{figure}[htbp]
  \centering
  \includegraphics[width=0.5\textwidth]{example-image}
  \caption{Overview of the design. Captions go below figures.}
  \label{fig:overview}
\end{figure}

\begin{figure}[htbp]
  \centering
  \begin{subfigure}[b]{0.3\textwidth}
    \includegraphics[width=\textwidth]{example-image-a}
    \caption{Loaded}
    \label{fig:steps-a}
  \end{subfigure}
  \hfill
  \begin{subfigure}[b]{0.3\textwidth}
    \includegraphics[width=\textwidth]{example-image-b}
    \caption{Released}
    \label{fig:steps-b}
  \end{subfigure}
  \hfill
  \begin{subfigure}[b]{0.3\textwidth}
    \includegraphics[width=\textwidth]{example-image-c}
    \caption{Launched}
    \label{fig:steps-c}
  \end{subfigure}
  \caption{The three stages of a launch.}
  \label{fig:steps}
\end{figure}

% ------------------------------------------------------------------------------
\section{Experimental setup}
\label{sec:setup}
% Why each experiment is done (to prove a claim), then how. Enough detail to
% reproduce it.
Each experiment tests one claim from \cref{sec:introduction}.
\Cref{tab:setup} lists the settings.

\begin{table}[htbp]
  \centering
  \caption{Settings of the experiments. Captions go above tables.}
  \label{tab:setup}
  \rowcolors{2}{Tue-accent!10}{white}
  \begin{tabular}{lccc}
    \toprule
    \color{Tue-accent}\textbf{Run} & \color{Tue-accent}\textbf{Mass (\si{\gram})}
    & \color{Tue-accent}\textbf{Angle (\si{\degree})}
    & \color{Tue-accent}\textbf{Repeats} \\
    \midrule
    A & 20 & 30 & 5 \\
    B & 20 & 45 & 5 \\
    C & 25 & 45 & \NA \\
    \bottomrule
  \end{tabular}
\end{table}

% ------------------------------------------------------------------------------
\section{Results}
\label{sec:results}
% Explain the results (do not only describe them) and show that they prove the
% claims.
% Graphs: axes with quantity and unit, the same ranges and legend order across
% graphs, readable in black and white, vector format (PDF) where possible.
Run B reached the longest distance, which confirms the claim that \todo{claim}.

% ------------------------------------------------------------------------------
\section{Conclusions}
\label{sec:conclusions}
% Restate the problem and the solution, and that the results prove the claims.
% No citations.
Restate the problem, the solution, and the results that prove the claims.

% ------------------------------------------------------------------------------
\printbibliography[heading=bibintoc, title={References}]

\clearpage
\appendix
\section{Calculations}
\label{app:calculations}
Long derivations and raw data go in appendices, referred to from the text as
\cref{app:calculations}.

\end{document}
